\section*{Prompt: Extract Relationships (Pass 2)}
\begin{description}
\item[Pipeline step:] Extraction --- Phase 2
\item[Models:] Same 3 models
\item[Source:] \texttt{misc/prompts.py} $\rightarrow$ \texttt{EXTRACT\_RELATIONSHIPS\_PROMPT}
\item[Called from:] \texttt{scripts/multi\_model\_mv.py} $\rightarrow$ \texttt{scripts/batch\_majority\_vote.py:build\_relationship\_jsonl()}
\end{description}
\begin{tcolorbox}[colback=yellow!8,colframe=yellow!50!black,left=4pt,right=4pt,top=4pt,bottom=4pt,arc=2pt]
\textbf{Note:} The \texttt{\{entities\_json\}} placeholder is replaced at runtime with the entities extracted in Phase 1.
\end{tcolorbox}

\begin{tcolorbox}[breakable,enhanced,colback=blue!3,colframe=blue!40!black,title={Prompt},fonttitle=\bfseries,left=8pt,right=8pt,top=6pt,bottom=6pt]

Your task is to act as a research assistant. You are given an academic paper and a set of \textbf{already-extracted entities} (constructs, measurement instruments, behaviors). Extract \textbf{relationships} between these entities that are explicitly supported by the paper text.

\subsection*{Already-Extracted Entities}

\{entities\_json\}

\subsection*{Relationship Extraction Rules}

\begin{itemize}
\item \textbf{ONLY use entities from the list above} as source and target. Do NOT introduce new entity names.
\item **\texttt{type} must be exactly one of:**
\begin{itemize}
\item \texttt{causes} --- explicit or strong causal claim ("X improves Y", "X enables Y", "X leads to Y")
\item \texttt{correlates} --- empirical association, regression, or "predicts" in the statistical sense ("X is associated with Y", "higher X corresponds to higher Y")
\item \texttt{measured\_by} --- a construct or behavior is evaluated/operationalized by a measurement instrument ("We evaluate X using benchmark Y", "We conducted a user study to assess X")
\item \texttt{related} --- a stated conceptual or empirical link that is NOT merely co-occurrence. The paper must describe HOW or WHY the two entities are connected.
\end{itemize}
\end{itemize}

\subsubsection*{Evidence Quality Rules (STRICT)}

\begin{itemize}
\item **\texttt{evidence}:\textbf{ must be a }quoted substring\textbf{ from the paper that }directly describes the connection** between source and target --- not just mentions both in the same sentence.
\item \textbf{Do NOT extract a relationship just because two entities appear in the same list, table, or sentence.} Listing instruments together (e.g. "we evaluate on MMLU, HumanEval, and GSM8K") does NOT create relationships between those instruments.
\item \textbf{Do NOT extract a relationship if the evidence only shows co-occurrence} (e.g. "We test reasoning and code generation" does NOT mean reasoning $\rightarrow$ code generation).
\item The evidence must answer: \textbf{"What does the paper say about how/why these two entities are connected?"} If you cannot answer that from the quote, do not include the relationship.
\end{itemize}

WRONG --- co-occurrence, no stated connection:
  source: "Reasoning", target: "Code generation", evidence: "We conduct experiments on reasoning tasks (language understanding, code generation, and mathematical reasoning)"

RIGHT --- stated connection:
  source: "Reasoning", target: "Aesthetic understanding", evidence: "to make further development of aesthetic understanding, reasoning is essential"

RIGHT --- measurement instrument relationship:
  source: "Bias", target: "PAIRS", type: "measured\_by", evidence: "We use PAIRS to probe social biases in VLMs"

RIGHT --- human evaluation relationship:
  source: "Helpfulness", target: "User preference study", type: "measured\_by", evidence: "We conducted a user preference study with 200 participants to evaluate helpfulness"

\begin{itemize}
\item \textbf{Allowed source--target patterns} (any combination justified by the text):
\begin{itemize}
\item construct $\leftrightarrow$ construct
\item construct $\leftrightarrow$ behavior
\item construct $\rightarrow$ measurement instrument (measured\_by)
\item measurement instrument $\rightarrow$ behavior
\item measurement instrument $\leftrightarrow$ measurement instrument (e.g. a benchmark suite contains a sub-benchmark, or two instruments measure overlapping domains)
\item behavior $\leftrightarrow$ behavior
\end{itemize}
\end{itemize}

\subsection*{Self-Check}

After drafting, verify:
\begin{itemize}
\item Every source and target name \textbf{exactly matches} an entity name from the list above.
\item Every evidence quote \textbf{appears} in the supplied paper text.
\item The evidence \textbf{describes a connection}, not just co-occurrence.
\item Direction and \texttt{type} match the paper's claims.
\item Prefer fewer, higher-quality relationships over many weak ones.
\end{itemize}

\subsection*{Output Format (JSON)}

Return a single JSON object with exactly:
\{\{
  "relationships": [
    \{\{
      "source": "Entity Name (from above list)",
      "target": "Entity Name (from above list)",
      "type": "causes | correlates | measured\_by | related",
      "evidence": "Quoted text from the paper that directly describes the connection (Section/Table ref)."
    \}\}
  ]
\}\}

If no relationships are supported, return \texttt{\{\{"relationships": []\}\}}.

Return only the JSON object, no additional text.
\end{tcolorbox}
